% Generado por sistema/backend ReporteTex.java -- no editar a mano.
% Datos: CSV importado (titulados) / CSV importado (empleadores)
\begin{itemize}
  \item \textbf{Becas y alianzas de financiamiento:} el 91.8\% de quienes cursaron un programa lo financió con recursos propios, y las pocas becas se concentran en ``Maestría o superior'' (17.2\% con beca, C7). Buscar financiamiento a través de organismos internacionales y convenios con instituciones y empresas para ampliar las becas (HU3).
  \item \textbf{Posgrados modulares} en las áreas de mayor interés (``IA'', ``Ing. de Software'', ``Redes y Seguridad''), para captar al 58.5\% de interesados que hoy elegiría otra institución (criterio 2.4.1), orientados por grupo de edad (C8): 26 o menos: Software y Base de Datos (52.9\%); 27–34: IA y Robótica (39.5\%); 35 o más: IA y Robótica (62.7\%); y rutas de continuidad (diplomado, especialidad, maestría) para quienes ya cursaron un programa (C11).
  \item \textbf{Formación según la etapa profesional}: especializaciones técnicas para los titulados junior y formación en gestión y dirección para los consolidados, porque la experiencia influye en la remuneración (C9) y en el cargo (C10).
  \item \textbf{Inserción laboral} (C12): reforzar la bolsa de trabajo y la vinculación con empleadores para el grupo de ``26 o menos'', donde el 41.0\% no trabaja.
  \item \textbf{Consulta periódica a empleadores}: incorporar a las organizaciones en los aspectos con menor acuerdo del criterio 2.4.2, en especial ``La Carrera consultó a los empleadores en su autoevaluación'' (35.8\%); ``La Carrera consultó a los empleadores sobre la pertinencia de la formación'' (43.4\%).
  \item \textbf{Actualización del plan de estudios}, que pide el 71.6\% de los titulados, con los contenidos que más mencionan los empleadores (``Inteligencia artificial y aprendizaje automático'', ``Ciberseguridad y seguridad de la información'', ``Análisis de datos, big data y BI'') (criterio 2.1.7).
  \item \textbf{Observatorio permanente de titulados}: mantener abierto el formulario de la herramienta y aprovechar el interés en la red de contactos (57.4\%) para aplicar la encuesta de forma periódica (criterio 3.2.2).
\end{itemize}
